% Source: https://artofproblemsolving.com/community/c3h1334418_mathcounts_hanook
% Stable identity: aops-topic-1334418-q01
Any common factor of $n^2-9$ and $n^2-4$ divides their difference $(n^2-4)-(n^2-9)=5$. So the GCF is $1$ or $5$, and it is $5$ exactly when $5\mid n^2-4=(n-2)(n+2)$, that is, when $n\equiv2$ or $n\equiv3\pmod5$ (then $5$ also divides $n^2-9=(n-3)(n+3)$).

From $1$ to $2015$ there are $403$ complete blocks of five consecutive integers, each containing one $n\equiv2$ and one $n\equiv3\pmod5$; and $2016\equiv1\pmod5$ does not count. So the answer is $2\cdot403=\boxed{806}$.

(For $n=2$ the fraction is undefined, but the two integers $n^2-9=-5$ and $n^2-4=0$ still have GCF $5$; if $n=2$ is left out because the fraction is undefined, the count is $805$.)
